%%%%%%%%%%%%%%%%
%%% lev 24 %%%
%%%%%%%%%%%%%%%%

\Chapter{24}

\Verse{1}
{
	\hOther{וַיְדַבֵּר יְהֹוָה אֶל מֹשֶׁה לֵּאמֹר׃}
}
{
	\eOther{24 And YHWH spoke to Moses, saying,}
}
{
}

\Verse{2}
{
	\hOther{צַו אֶת בְּנֵי יִשְׂרָאֵל וְיִקְחוּ אֵלֶיךָ שֶׁמֶן זַיִת זָךְ כָּתִית לַמָּאוֹר לְהַעֲלֹת נֵר תָּמִיד׃}
}
{
	\eOther{
		“Command the children of Israel that they shall take clear
		pressed olive oil to you for the light, to keep up a lamp
		always.
	}
}
{
}

\Verse{3}
{
	\hOther{מִחוּץ לְפָרֹכֶת הָעֵדֻת בְּאֹהֶל מוֹעֵד יַעֲרֹךְ אֹתוֹ אַהֲרֹן מֵעֶרֶב עַד בֹּקֶר לִפְנֵי יְהֹוָה תָּמִיד חֻקַּת עוֹלָם לְדֹרֹתֵיכֶם׃}
}
{
	\eOther{
		Outside of the pavilion of the Testimony, in the Tent of
		Meeting, Aaron shall arrange it from evening until morning
		in front of YHWH always, an eternal law through your
		generations.
	}
}
{
}

\Verse{4}
{
	\hOther{עַל הַמְּנֹרָה הַטְּהֹרָה יַעֲרֹךְ אֶת הַנֵּרוֹת לִפְנֵי יְהֹוָה תָּמִיד׃}
}
{
	\eOther{
		He shall arrange the lamps on the pure menorah in front of
		YHWH always.
	}
}
{
}

\Verse{5}
{
	\hOther{וְלָקַחְתָּ סֹלֶת וְאָפִיתָ אֹתָהּ שְׁתֵּים עֶשְׂרֵה חַלּוֹת שְׁנֵי עֶשְׂרֹנִים יִהְיֶה הַחַלָּה הָאֶחָת׃}
}
{
	\eOther{
		“And you shall take fine flour and bake it into twelve
		loaves: each loaf shall be two-tenths of a measure.
	}
}
{
}

\Verse{6}
{
	\hOther{וְשַׂמְתָּ אוֹתָם שְׁתַּיִם מַעֲרָכוֹת שֵׁשׁ הַמַּעֲרָכֶת עַל הַשֻּׁלְחָן הַטָּהֹר לִפְנֵי יְהֹוָה׃}
}
{
	\eOther{
		And you shall set them in two rows, six to a row, on the
		pure table in front of YHWH.
	}
}
{
}

\Verse{7}
{
	\hOther{וְנָתַתָּ עַל הַמַּעֲרֶכֶת לְבֹנָה זַכָּה וְהָיְתָה לַלֶּחֶם לְאַזְכָּרָה אִשֶּׁה לַיהֹוָה׃}
}
{
	\eOther{
		And you shall put clear frankincense on each row, and it
		will be a representative portion for the bread, an
		offering by fire to YHWH.
	}
}
{
}

\Verse{8}
{
	\hOther{בְּיוֹם הַשַּׁבָּת בְּיוֹם הַשַּׁבָּת יַעַרְכֶנּוּ לִפְנֵי יְהֹוָה תָּמִיד מֵאֵת בְּנֵי יִשְׂרָאֵל בְּרִית עוֹלָם׃}
}
{
	\eOther{
		On the Sabbath day, on the Sabbath day, he shall arrange
		it in front of YHWH always from the children of Israel, an
		eternal covenant.
	}
}
{
%	\footnote{
%		On the Sabbath day, on the Sabbath day. Meaning: on every Sabbath.
%	}
}

\Verse{9}
{
	\hOther{וְהָיְתָה לְאַהֲרֹן וּלְבָנָיו וַאֲכָלֻהוּ בְּמָקוֹם קָדֹשׁ כִּי קֹדֶשׁ קדָשִׁים הוּא לוֹ מֵאִשֵּׁי יְהֹוָה חק עוֹלָם׃}
}
{
	\eOther{
		And it shall be Aaron’s and his sons’, and they shall eat
		it in a holy place, because it is the holy of holies to
		him from YHWH’s offerings by fire, an eternal law.”
	}
}
{
}

\Verse{10}
{
	\hOther{וַיֵּצֵא בֶּן אִשָּׁה יִשְׂרְאֵלִית וְהוּא בֶּן אִישׁ מִצְרִי בְּתוֹךְ בְּנֵי יִשְׂרָאֵל וַיִּנָּצוּ בַּמַּחֲנֶה בֶּן הַיִּשְׂרְאֵלִית וְאִישׁ הַיִּשְׂרְאֵלִי׃}
}
{
	\eOther{
		And a son of an Israelite woman—and he was a son of an
		Egyptian man—went out among the children of Israel, and
		the son of the Israelite woman fought with an Israelite
		man in the camp. 24:10–23. a son of an Israelite woman....
		they battered him with stone. It is curious that this
		single story occurs here in a section fifteen chapters
		long that otherwise simply states the law without any
		accompanying narratives. This story’s significance may lie
		in the fact that it states that the man in question is the
		son of an Israelite woman and an Egyptian father, and the
		deity’s decision includes the notation that the law
		applies equally to an alien and to a citizen. It therefore
		appears that the offender in this story is an alien (see
		the comment on Lev 24:16,22), and the account therefore
		has the value not only of dramatizing the law concerning
		blasphemy, but also of dramatizing the principle of
		equality before the law: “You shall have one judgment: it
		will be the same for the alien and the citizen” (24:22).
		Indeed, the group of laws that were described above as
		reflecting the notion of equivalence of justice (“eye for
		eye … ”) is embedded in the middle of this very story. The
		story may appear here, therefore, because it illustrates
		several important laws and an extremely important legal
		principle. The story is also interesting in parallel with
		the only other narrative in Leviticus, the story of the
		consecration of the priesthood. Like that other story,
		which culminates in the deaths of Aaron’s sons, the
		account of the blasphemy expresses what is at stake in the
		people’s young relationship with their God. Further, in
		the earlier account, it is God who deals with the
		offenders, Nadab and Abihu, directly; whereas in this
		account of the blasphemy, humans themselves must deal
		justice. The two stories thus convey together the idea
		that the law is both a divine and a human concern.
	}
}
{
}

\Verse{11}
{
	\hOther{וַיִּקֹּב בֶּן הָאִשָּׁה הַיִּשְׂרְאֵלִית אֶת הַשֵּׁם וַיְקַלֵּל וַיָּבִיאוּ אֹתוֹ אֶל מֹשֶׁה וְשֵׁם אִמּוֹ שְׁלֹמִית בַּת דִּבְרִי לְמַטֵּה דָן׃}
}
{
	\eOther{
		And the son of the Israelite woman profaned the name and
		cursed. And they brought him to Moses. And his mother’s
		name was Shelomith, daughter of Dibri, of the tribe of
		Dan.
	}
}
{
}

\Verse{12}
{
	\hOther{וַיַּנִּיחֻהוּ בַּמִּשְׁמָר לִפְרֹשׁ לָהֶם עַל פִּי יְהֹוָה׃}
}
{
	\eOther{
		And they left him under watch, to determine it for them by
		YHWH’s word.
	}
}
{
}

\Verse{13}
{
	\hOther{וַיְדַבֵּר יְהֹוָה אֶל מֹשֶׁה לֵּאמֹר׃}
}
{
	\eOther{And YHWH spoke to Moses, saying,}
}
{
}

\Verse{14}
{
	\hOther{הוֹצֵא אֶת הַמְקַלֵּל אֶל מִחוּץ לַמַּחֲנֶה וְסָמְכוּ כל הַשֹּׁמְעִים אֶת יְדֵיהֶם עַל רֹאשׁוֹ וְרָגְמוּ אֹתוֹ כּל הָעֵדָה׃}
}
{
	\eOther{
		“Take out the one who cursed to the outside of the camp,
		and let all who heard lay their hands on his head, and all
		the congregation shall batter him.
	}
}
{
%	\footnote{
%		lay their hands on his head. The community’s elders lay their hands on
%		the head of a bull that is slaughtered as a sacrifice for a sin (4:15);
%		and Aaron and his sons lay their hands on the head of the animal that is
%		sacrificed as a sin offering (Lev 8:14). And so the mention of this
%		ritual here brings up a frightful image of what is to come for the man
%		who cursed, and it emphasizes the unusual solemnity of this execution
%		for the crime of cursing God. Footnote 24:14. batter. Meaning: stone to
%		death.
%	}
}

\Verse{15}
{
	\hOther{וְאֶל בְּנֵי יִשְׂרָאֵל תְּדַבֵּר לֵאמֹר אִישׁ אִישׁ כִּי יְקַלֵּל אֱלֹהָיו וְנָשָׂא חֶטְאוֹ׃}
}
{
	\eOther{
		“And you shall speak to the children of Israel, saying:
		Any man who curses his God: he shall bear his sin.
	}
}
{
}

\Verse{16}
{
	\hOther{וְנֹקֵב שֵׁם יְהֹוָה מוֹת יוּמָת רָגוֹם יִרְגְּמוּ בוֹ כּל הָעֵדָה כַּגֵּר כָּאֶזְרָח בְּנקְבוֹ שֵׁם יוּמָת׃}
}
{
	\eOther{
		And one who profanes YHWH’s name shall be put to death.
		All the congregation shall batter him. The same for the
		alien and the citizen: when he profanes the name he shall
		be put to death.
	}
}
{
%	\footnote{
%		the same for the alien and the citizen. This phrase occurs at both the
%		beginning and the end of this section on equivalence in justice. This
%		principle of treating an alien like a citizen comes up often in the
%		Torah, but its double emphasis here in the story of a man who has an
%		Israelite mother and an Egyptian father is significant. It indicates
%		that this man is an alien, because there would be no reason to make a
%		special point of this principle if he were a citizen (that is, a native
%		Israelite). It is further evidence that in biblical times a person’s
%		identification was based on the father (patrilineal descent). In every
%		case we have encountered thus far, the woman and children take the
%		ethnic and religious identity of the woman’s husband—Sarah, Rebekah,
%		Leah, Rachel, Zipporah, and their children. And all sons of Levite
%		fathers are Levites, regardless of the tribal origins of their mothers.
%		Likewise, all tribal identification is through the father (see the
%		comment on Num 36:3). Identification in the Torah is patrilineal. This
%		continues through the Tanak until the time of Ezra and Nehemiah. At that
%		point it appears to change, because when Nehemiah opposes marriages of
%		Jewish men to women from certain other groups, he instructs that their
%		children be excluded (Nehemiah 13). And in postbiblical Judaism,
%		identification came to be based on the mother (matrilineal descent). In
%		the Torah, however, it follows the father in every case. Identification
%		by patrilineal descent (in addition to matrilineal descent, not in place
%		of it) is at present a matter of dispute among the movements in Judaism.
%	}
}

\Verse{17}
{
	\hOther{וְאִישׁ כִּי יַכֶּה כּל נֶפֶשׁ אָדָם מוֹת יוּמָת׃}
}
{
	\eOther{
		“And a man who will strike any human’s life shall be put
		to death,
	}
}
{
%	\footnote{
%		strike any human’s life. Meaning: strike to death.
%	}
}

\Verse{18}
{
	\hOther{וּמַכֵּה נֶפֶשׁ בְּהֵמָה יְשַׁלְּמֶנָּה נֶפֶשׁ תַּחַת נָפֶשׁ׃}
}
{
	\eOther{
		and one who strikes an animal’s life shall pay for it: a
		life for a life.
	}
}
{
}

\Verse{19}
{
	\hOther{וְאִישׁ כִּי יִתֵּן מוּם בַּעֲמִיתוֹ כַּאֲשֶׁר עָשָׂה כֵּן יֵעָשֶׂה לּוֹ׃}
}
{
	\eOther{
		And a man who will make an injury in his fellow: as he has
		done, so it shall be done to him.
	}
}
{
}

\Verse{20}
{
	\hOther{שֶׁבֶר תַּחַת שֶׁבֶר עַיִן תַּחַת עַיִן שֵׁן תַּחַת שֵׁן כַּאֲשֶׁר יִתֵּן מוּם בָּאָדָם כֵּן יִנָּתֶן בּוֹ׃}
}
{
	\eOther{
		A break for a break, an eye for an eye, a tooth for a
		tooth: as he will make an injury in a human, so it shall
		be made in him.
	}
}
{
%	\footnote{
%		an eye for an eye. Perhaps the most perplexing of the ethical laws is
%		the principle of justice expressed in the formulation “an eye for an
%		eye.” It has frequently been cited as evidence of the stern character of
%		YHWH, but that is a misunderstanding. In its context in Leviticus it
%		applies solely to human justice. YHWH Himself frequently follows a more
%		relenting course than that, from the golden calf event to a series of
%		reprieves for seemingly undeserving individuals and communities in
%		subsequent books of the Tanak. As for the meaning of this formulation
%		for human justice, we must read it in its context, where the basic
%		principle appears to be that punishment should correspond to the crime
%		and never exceed it: A man who will make an injury in his fellow: as he
%		has done, so it shall be done to him. A break for a break, an eye for an
%		eye, a tooth for a tooth: as he will make an injury in a human, so it
%		shall be made in him. And one who strikes [mortally] an animal shall pay
%		for it, and one who strikes [mortally] a human shall be put to death.
%		You shall have one judgment: it will be the same for the alien and the
%		citizen. The common principle of all the instances in this group appears
%		to be one of equivalence in justice. To that extent it is a prohibition
%		against imposing punishments that do not suit the crime; e.g., execution
%		as a punishment for killing an animal or for wounding a man; or more
%		severe punishments for lower-class defendants or foreigners than for
%		others. The question remains as to whether ”a break for a break, an eye
%		for an eye” was meant literally or if it likewise referred figuratively
%		to equivalence of punishment: not requiring the court actually to gouge
%		the eye of a man convicted of gouging, but only to exact an appropriate
%		monetary compensation for the victim. The text certainly appears to be
%		speaking quite literally. It is set in the middle of a story of an
%		actual stoning, and the wording would not ordinarily be taken as
%		figurative were it not for our surprise at the severity of it in
%		comparison to our own current laws. Still, we must be cautious about our
%		ability to judge the meaning of idiomatic expressions in the Hebrew of
%		an ancient age, a Hebrew in which, for example, the expression “to know
%		God face-to-face” explicitly did not mean actually to see the deity’s
%		face (Exod 33:11; cf. 33:20,23). Further, the earliest postbiblical
%		Jewish sources already understood “an eye for an eye” to mean monetary,
%		and not literal, compensation. In the present state of our knowledge,
%		therefore, we may be limited to saying that this text either (1) meant
%		equivalent monetary compensation, or else (2) initially meant literal
%		physical punishment but was reinterpreted as monetary equivalence at a
%		relatively early stage in legal history (in the early postbiblical
%		period at the latest).
%	}
}

\Verse{21}
{
	\hOther{וּמַכֵּה בְהֵמָה יְשַׁלְּמֶנָּה וּמַכֵּה אָדָם יוּמָת׃}
}
{
	\eOther{
		And one who strikes an animal shall pay for it, and one
		who strikes a human shall be put to death.
	}
}
{
}

\Verse{22}
{
	\hOther{מִשְׁפַּט אֶחָד יִהְיֶה לָכֶם כַּגֵּר כָּאֶזְרָח יִהְיֶה כִּי אֲנִי יְהֹוָה אֱלֹהֵיכֶם׃}
}
{
	\eOther{
		You shall have one judgment: it will be the same for the
		alien and the citizen; because I am YHWH, your God.”
	}
}
{
}

\Verse{23}
{
	\hOther{וַיְדַבֵּר מֹשֶׁה אֶל בְּנֵי יִשְׂרָאֵל וַיּוֹצִיאוּ אֶת הַמְקַלֵּל אֶל מִחוּץ לַמַּחֲנֶה וַיִּרְגְּמוּ אֹתוֹ אָבֶן וּבְנֵי יִשְׂרָאֵל עָשׂוּ כַּאֲשֶׁר צִוָּה יְהֹוָה אֶת מֹשֶׁה׃}
}
{
	\eOther{
		And Moses spoke to the children of Israel. And they
		brought out the one who cursed to the outside of the camp,
		and they battered him with stone. And the children of
		Israel had done as YHWH commanded Moses.
	}
}
{
}
